\documentclass[11pt]{article}
\usepackage[a4paper,margin=1in]{geometry}
\usepackage{amsmath,amssymb,amsthm,mathtools}
\usepackage[hidelinks]{hyperref}
\usepackage{booktabs}

\newtheorem{theorem}{Theorem}
\newtheorem{proposition}[theorem]{Proposition}
\newtheorem{lemma}[theorem]{Lemma}
\newcommand{\Cl}{\operatorname{Cl}}
\newcommand{\Tr}{\operatorname{Tr}}
\newcommand{\Q}{\mathbb Q}
\newcommand{\Z}{\mathbb Z}

\title{Twisted-Convolution Identities in Dimensions Seven and Eight:\\
Shintani Height Rigidity and CM Descent}
\author{Hainan Zhao\\
\small Independent Researcher\\
\small\texttt{hainzhao@gmail.com}}
\date{28 July 2026}

\begin{document}
\maketitle

\begin{abstract}
We prove the formal Twisted Convolution Conjecture of
Appleby--Flammia--Kopp unconditionally for the discriminant-$32$
dimension-seven stratum and for every admissible tuple in dimension
eight.  In dimension seven, conductor lowering places every characteristic in
one of six maximal-order ray groups.  Shintani's proved index-two
algebraicity theorem gives a safe exponent $16128$; rigorous Arb
enclosures, Voutier's height gap, exact Artin labels, and a
degree-$48$ compositum certificate identify the complete signed
overlap packet and make all $441$ rank-two minors vanish for both
formal shifts.  Dimension eight has two admissible form conductors.
For conductor three, genuine linear reinduction to quartic ray
characters over $\Q(\sqrt{-6})$ and Stark's proved
imaginary-quadratic rank-one theorem determine the two formerly
missing orientations.  For conductor one, the relevant ray characters
are quadratic; analytic class-number formulas determine their units,
and an exact phase calculation in $\Q(\sqrt5)$ orients all $63$
nonzero cocycle values.  Exact quotient-ring arithmetic makes all
$784$ minors vanish for both shifts in each dimension-eight stratum.
No unproved Stark conjecture or numerical recognition is used as a
logical input. Dimension-seven admissibility also permits discriminant
$8$; that second stratum is isolated explicitly but is not claimed
here.
\end{abstract}

\section{Formal TCC and conventions}

We use the conventions of Appleby--Flammia--Kopp (AFK)~\cite{AFK}.
An admissible tuple $t=(d,r,Q)$ contains a primitive indefinite
integral binary quadratic form $Q$, a positive real fixed point
$\rho_t$, a level-$d$ stabilizer $A_t$, and the matrix $L_{z,t}$.
For $\boldsymbol p\in\Z^2$, let $\mathcal I_{\boldsymbol p}$ be a
complete set of representatives modulo $d$ containing
$\boldsymbol0$ and $\boldsymbol p$.  A residue
$\lambda\in\Z/d\Z$ is a \emph{shift} when its AFK coprimality
condition holds and
\[
 \sum_{\boldsymbol q\in\mathcal I_{\boldsymbol p}}
 \omega_d^{\,r\langle\boldsymbol p,
  (\lambda I+L_{z,t})\boldsymbol q\rangle}
 \operatorname{shin}_{A_t}^{\boldsymbol q/d}(\rho_t)
 \operatorname{shin}_{A_t^{-1}}^{
       (\boldsymbol q-\boldsymbol p)/d}(\rho_t)
 =d^2\delta_{\boldsymbol p,\boldsymbol0}^{(d)}
\tag{TCC}
\]
for every $\boldsymbol p$.  Write $\mathcal Z_t$ for the shift set.
The formal TCC asserts $0,1\in\mathcal Z_t$ and equality of shift
sets for admissible forms of the same discriminant.

Our reciprocal double sine is
\[
 S_2(z\mid\omega_1,\omega_2)
 :=\frac{\Gamma_2(z\mid\omega_1,\omega_2)}
 {\Gamma_2(\omega_1+\omega_2-z\mid\omega_1,\omega_2)}.
\tag{1}
\]
It is the reciprocal of Kopp's
$\operatorname{Sin}_{2,\mathrm{Kopp}}$ convention
\cite[Definition~4.22]{Kopp}.  In particular,
\[
 S_2(z+1\mid\beta,1)
 =2\sin(\pi z/\beta)S_2(z\mid\beta,1).
\tag{2}
\]
The place labels $\infty_1,\infty_2$ always refer to the positive and
negative square-root embeddings printed in the certificates.  PARI
orders the negative root first in the quadratic fields used below;
all raw PARI conductor vectors are translated to the displayed place
labels before use.

\begin{proposition}[Finite reconstruction bridge]\label{prop:bridge}
Suppose the convention-matched AFK values
$\nu_{\boldsymbol p}$ are inserted into the normalized Weyl
reconstruction for a permitted twist.  Then the reconstructed matrix
$\Pi$ satisfies $\Tr\Pi=1$, and equation \textup{(TCC)} is equivalent
to $\Pi^2=\Pi$.  If all rank-two minors vanish and $\Pi\ne0$, then
$\Pi$ is a rank-one idempotent and the corresponding residue is a
shift.
\end{proposition}

\begin{proof}
Expand the Weyl coefficient of $\Pi^2-\Pi$, use the Weyl multiplication
law, and reindex by the displacement difference.  The AFK endpoint
correction changes the raw diagonal value
$\sqrt{d+1}$ to the normalized zero overlap one.  The resulting
coefficient is exactly the left side of \textup{(TCC)} minus its
$d^2\delta$ term.  Fourier inversion is injective.  Finally, vanishing
of the $2\times2$ minors gives rank at most one, while trace one gives
rank exactly one and idempotency.
\end{proof}

\begin{theorem}\label{thm:main}
For every admissible dimension-seven tuple $t$ whose form has
discriminant $32$, and for every admissible dimension-eight tuple $t$,
\[
 0,1\in\mathcal Z_t.
\]
If $t,t'$ have forms of the same discriminant, then
$\mathcal Z_t=\mathcal Z_{t'}$.
\end{theorem}

\section{Dimension seven}\label{sec:d7}

\subsection{Admissible strata and scope}

The admissibility recurrence forces
\[
 d=7\Longrightarrow
 r=1,\quad n=8,\quad K_7=\Q(\sqrt2),\quad
 j=m=1,\quad f_1=2.
\]
The form conductor may therefore be $f=1$ or $f=2$.  Both strata
exist; convenient representatives are
\[
 Q_{7,1}=\langle1,-4,2\rangle,\quad
 \operatorname{disc}Q_{7,1}=8,
 \qquad
 Q_{7,2}=\langle1,-6,1\rangle,\quad
 \operatorname{disc}Q_{7,2}=32.
\]
This section proves the formal TCC for $Q_{7,2}$ and transports it
within discriminant $32$.  It does not use form covariance to cross
from discriminant $32$ to discriminant $8$, and it makes no claim for
$Q_{7,1}$.

\subsection{Conductor lowering and ray classes}

Put
\[
 K_7=\Q(\sqrt2),\qquad
 \alpha=2+\sqrt2,\qquad
 \beta_7=3+2\sqrt2.
\]
Here $\beta_7$ is the positive fixed point of $Q_{7,2}$, whereas
$\alpha$ is the reduced maximal-order point used only after conductor
lowering.  With
\[
 L_7=\begin{pmatrix}6&-1\\1&0\end{pmatrix},\qquad
 A_{7,2}=L_7^3=\begin{pmatrix}204&-35\\35&-6\end{pmatrix},
\]
the exact Rademacher calculation gives
\[
 \Phi_{\boldsymbol p}
 =\zeta_{56}^{\,7-32Q_{7,2}(\boldsymbol p)}.
\]
The determinant-two map
\[
 B=\begin{pmatrix}2&-1\\0&1\end{pmatrix}
\]
sends $\alpha$ to $\beta_7$.  For a characteristic $(a,b)/7$, its two
preimages are
\[
 s_j=\frac1{14}\binom{a+b+7j}{2b},\qquad j=0,1.
\tag{3}
\]
Kopp's inverse ray map assigns to $s_j$ the principal element
\[
 \gamma_{a,b,j}=2b\sqrt2+3b-a-7j.
\tag{4}
\]
Removing the common ideal divisor of $(14)$ and $(\gamma)$ produces
exactly six lowered moduli.  Their one-place ray groups and Kopp
exponents are
\[
\begin{array}{c|c|c}
\text{modulus HNF}&\text{ray group}&n\\ \hline
\left(\begin{smallmatrix}14&0\\0&14\end{smallmatrix}\right)&C_6\times C_2&1\\
\left(\begin{smallmatrix}7&0\\0&7\end{smallmatrix}\right)&C_6&1\\
\left(\begin{smallmatrix}14&6\\0&2\end{smallmatrix}\right)&C_2&1\\
\left(\begin{smallmatrix}14&8\\0&2\end{smallmatrix}\right)&C_2&1\\
\left(\begin{smallmatrix}7&3\\0&1\end{smallmatrix}\right)&C_2&2\\
\left(\begin{smallmatrix}7&4\\0&1\end{smallmatrix}\right)&1&2.
\end{array}
\tag{5}
\]
The common determinant-one stabilizer of the lowered factors at
$\alpha$ is
\[
 A_{\rm low}=\begin{pmatrix}239&-140\\70&-41\end{pmatrix}.
\tag{6}
\]
The exact inverse-ray, gcd, sign-class, and stabilizer computations are
printed for all $96$ lifted factors by
\path{scripts/explore_dimension_seven_conductor_lowering.gp}.

\subsection{Powered algebraicity and packet identification}

The full one-place ray-$14$ field has relative group
$C_6\times C_2$.  In its normal closure, the maximal absolutely
abelian subfield has index two and the applicable imaginary-quadratic
base is $\Q(\sqrt{-7})$.  Auditing all $32$ divisors of its conductor
gives
\[
 m_{\mathfrak d}=
 \begin{cases}
 12h_{\Q(\sqrt{-7})}n(S_{\mathfrak d}),&\mathfrak d=(1),\\
 12f_{\mathfrak d}n(S_{\mathfrak d}),&\mathfrak d\ne(1),
 \end{cases}
 \qquad
 m=\operatorname{lcm}_{\mathfrak d\mid\mathfrak c}
 m_{\mathfrak d}=16128.
\tag{7}
\]
Write
$\mathfrak c=\mathfrak p_2^3\bar{\mathfrak p}_2^3\mathfrak p_7$ and
$\mathfrak d(e_1,e_2,e_7)
=\mathfrak p_2^{e_1}\bar{\mathfrak p}_2^{e_2}
\mathfrak p_7^{e_7}$. The complete divisor audit is below.
Each triple records
\((|H_{\mathfrak d}|,n(S_{\mathfrak d}),m_{\mathfrak d})\);
here $|H_{\mathfrak d}|$ is the ray-group order, and $w_0/w_1$
records $w(\mathfrak d)$ for $e_7=0/1$. The two triple columns have
$e_7=0$ and $e_7=1$, respectively.
\[
\begin{array}{cc|c|c|c}
e_1&e_2&w_0/w_1&e_7=0&e_7=1\\ \hline
0&0&2/1&(1,96,1152)&(3,16,1344)\\
0&1&2/1&(1,96,2304)&(3,16,2688)\\
0&2&1/1&(1,48,2304)&(6,8,2688)\\
0&3&1/1&(2,24,2304)&(12,4,2688)\\
1&0&2/1&(1,96,2304)&(3,16,2688)\\
1&1&2/1&(1,96,2304)&(3,16,2688)\\
1&2&1/1&(1,48,2304)&(6,8,2688)\\
1&3&1/1&(2,24,2304)&(12,4,2688)\\
2&0&1/1&(1,48,2304)&(6,8,2688)\\
2&1&1/1&(1,48,2304)&(6,8,2688)\\
2&2&1/1&(2,24,1152)&(12,4,1344)\\
2&3&1/1&(4,12,1152)&(24,2,1344)\\
3&0&1/1&(2,24,2304)&(12,4,2688)\\
3&1&1/1&(2,24,2304)&(12,4,2688)\\
3&2&1/1&(4,12,1152)&(24,2,1344)\\
3&3&1/1&(8,6,576)&(48,1,672)
\end{array}
\]
Thus the lcm of all thirty-two displayed exponents is $16128$.
The full HNF ideals and independently recomputed ray orders are in the
archived transcript
\[
\texttt{certificates/dimension-seven-shintani-divisors.txt}.
\]
Here the unit-congruence factor
$w(\mathfrak d)$ is computed separately for every divisor; four
divisors, including the trivial one, have $w=2$.  Shintani's proved
theorem~\cite{Shintani1978} therefore makes the powered packet
algebraic.

The sixteen squared-overlap representatives satisfy the reciprocal
polynomial
\[
\begin{aligned}
P(X)={}&X^{16}-(37+28\sqrt2)X^{15}
 +(1212+854\sqrt2)X^{14}\\
&-(20685+14630\sqrt2)X^{13}
 +(210371+148750\sqrt2)X^{12}\\
&-(1313872+929054\sqrt2)X^{11}
 +(5031845+3558044\sqrt2)X^{10}\\
&-(11531249+8153832\sqrt2)X^9
 +(15288774+10810788\sqrt2)X^8\\
&-(11531249+8153832\sqrt2)X^7
 +(5031845+3558044\sqrt2)X^6\\
&-(1313872+929054\sqrt2)X^5
 +(210371+148750\sqrt2)X^4\\
&-(20685+14630\sqrt2)X^3
 +(1212+854\sqrt2)X^2\\
&-(37+28\sqrt2)X+1.
\end{aligned}
\]
It factors over $\Q(\sqrt2)$ with degrees $2,2,12$. Exact
class-field comparisons identify those factors with the two quadratic
lowered fields $P_+,P_-$ and the one-place ray-$14$ field. The
following rational intervals each contain exactly one real root by
Sturm's theorem; the last column is the exact ray log or quadratic
class label.
\[
\begin{array}{c|c|c@{\quad}c|c|c}
(a,b)&\text{interval}&\text{label}&(a,b)&\text{interval}&\text{label}\\ \hline
(2,2)&(.043513,.043514)&[2,1]&(1,4)&(.079508,.079509)&P_+:[1,1]\\
(1,3)&(.080124,.080125)&[3,0]&(1,2)&(.104688,.104689)&P_-:[1]\\
(1,1)&(.146404,.146405)&[4,1]&(0,6)&(.169377,.169378)&[3,1]\\
(0,5)&(.315227,.315228)&[1,1]&(0,4)&(.671533,.671534)&[2,0]\\
(0,3)&(1.489129,1.489130)&[5,1]&(0,2)&(3.172313,3.172314)&[4,0]\\
(0,1)&(5.903986,5.903987)&[0,0]&(2,6)&(6.830385,6.830386)&[1,0]\\
(3,6)&(9.552165,9.552166)&P_-:[0]&(4,4)&(12.480646,12.480647)&[0,1]\\
(3,5)&(12.577346,12.577347)&P_+:[0,0]&(4,5)&(22.981630,22.981631)&[5,0]
\end{array}
\]
The split primes above $17$ and $41$ have ray logs $[1,0]$ and
$[0,1]$ and act by automorphisms $5$ and $10$, respectively. Local
Frobenius congruences certify these actions and every table label in
\path{scripts/dimension_seven_artin_labels.gp}.

Rigorous Arb evaluation of the eight independent analytic values gives
the raw logarithmic enclosure
\[
 \epsilon_{\log}\le 5.86\cdot10^{-11},\qquad
 16128\epsilon_{\log}\le9.45\cdot10^{-7}.
\]
The complete balls are archived in
\[
\texttt{certificates/dimension-seven-double-sine-intervals.txt}.
\]
Voutier's lower bound~\cite{Voutier1996} is greater than
$5.22\cdot10^{-5}$ for every relevant degree. The quotient of an
analytic value by its assigned algebraic root must therefore be a root
of unity; positivity and the exact cocycle multipliers select the
displayed root.

\subsection{Exact finite certificate}

Let $H$ be the degree-$24$ signed overlap field and
$C=\Q(\zeta_{56})$.  Their intersection has degree $12$.  The standard
component is selected by the exact identity
\[
 c_H=\zeta_{56}+\zeta_{56}^{-1},\qquad
 1.98<c_H<1.99,
\tag{8}
\]
and it also satisfies
\[
 \sqrt2=\zeta_{56}^{7}+\zeta_{56}^{-7}.
\tag{9}
\]
In the resulting degree-$48$ compositum,
\path{scripts/dimension_seven_exact_tcc.gp} verifies for each formal
shift
\[
 \Tr\Pi=1,\qquad \Pi^2=\Pi,\qquad
 441\text{ rank-two minors vanish}.
\tag{10}
\]
Proposition~\ref{prop:bridge} proves both shifts for the
discriminant-$32$ tuple $Q_{7,2}$. The order of discriminant $32$ has
wide class number one. Unconditional $\mathrm{GL}_2(\Z)$ covariance
therefore transports the result to every admissible dimension-seven
form of discriminant $32$, and no farther. The discriminant-$8$
stratum requires an independent analytic packet and remains open.

\section{Dimension eight}\label{sec:d8}

Dimension-eight admissibility has $j=2$ and $f_2=3$.  The form
conductor is therefore one or three, giving discriminants $5$ and
$45$.  These strata must be proved independently.

\subsection{Conductor three: linear CM reinduction}

For the canonical discriminant-$45$ form, the order-ray group is
identified with the maximal-order ray group
\[
 \Cl_{(8)\infty_2}(\Z+3\mathcal O_{\Q(\sqrt5)})
 \simeq
 \Cl_{(24)\infty_2}(\mathcal O_{\Q(\sqrt5)})
 \simeq C_4\times C_2^2.
\tag{11}
\]
Kopp's conductor-lowering formula supplies the exact three-factor
dictionary.  The fifteen lower characteristics lie in an index-two
ray-$12$ field; Shintani exponent $576$, Arb isolation, and height
rigidity prove that entire lower packet.

The primitive packet contains two conjugate pairs of quartic
characters.  On the full two-place ray group, base conjugation does
not literally send either one-place character to its inverse because
it swaps the labeled infinite place.  The projective quotient is
nevertheless the same order-two character in both cases, and the
associated biquadratic field has quadratic subfields
\[
 \Q(\sqrt5),\qquad\Q(\sqrt{-6}),\qquad\Q(\sqrt{-30}).
\tag{12}
\]
This projective observation is not used as a substitute for linear
equivalence.  The script
\path{scripts/dimension_eight_linear_cm_reinduction.gp} compares the
complete induced characters on the degree-$16$ normal closure and
proves genuine linear reinduction, without a scalar twist, from
quartic ray characters over $\Q(\sqrt{-6})$.  The alternative base
$\Q(\sqrt{-30})$ gives an independent check.

\begin{lemma}[Ray labels and local factors]\label{lem:cm-label}
Let
\[
 \mathfrak m_M
 =\mathfrak p_2^3\mathfrak p_3\mathfrak p_5,
\]
with no infinite component, where the prime ideals are the exact
factors printed by
\path{scripts/dimension_eight_cm_unit_lattice.gp}. The ray group
$\Cl_{\mathfrak m_M}(M)$ has cyclic coordinates $[8,4]$, and in those
coordinates the reinduced characters are
\[
 \psi_0=[6,1],\qquad \psi_1=[2,1].
\tag{13}
\]
They have the same complete Artin $L$-functions as the original
$\Q(\sqrt5)$ characters $[1,0,0]$ and $[1,1,0]$, respectively. With
\[
 S_M=\{\infty,\mathfrak p_2,\mathfrak p_3,\mathfrak p_5\},
\tag{14}
\]
every finite member is ramified for $\psi_b$, so
$L_{M,S_M}(s,\psi_b)=L_M(s,\psi_b)$.
\end{lemma}

\begin{proof}
The exact normal-closure calculation leaves precisely two faithful
ray characters for each kernel, an inverse pair.  For packet zero the
coefficient at $n=5$ is $-i$ for $[6,1]$ and $+i$ for $[2,3]$; the
source coefficient is $-i$.  For packet one it is $+i$ for $[2,1]$
and $-i$ for $[6,3]$; the source coefficient is $+i$.  This single
exact separating coefficient selects among an exhaustive
two-element set; it is not finite-coefficient recognition.
\path{scripts/dimension_eight_cm_unit_lattice.gp} fixes the relative
factors before constructing their ray groups and prints these four
coefficients.

The two source characters have full conductor $(24)\infty_2$.
The same audit proves that both CM conductors equal
$\mathfrak m_M$ and that $\psi_b$ is ramified at every
finite member of $S_M$.  Thus those local factors are already one.
Linear reinduction gives the equality of the complete functions
$L_M(s,\psi_b)=L_K(s,\chi_b)$; no unramified source Euler factor at
$5$ is removed. The archimedean parity is part of the same induced
character equality.
\end{proof}

\begin{lemma}[Normalized imaginary-quadratic Stark resolvent]
\label{lem:stark-normalization}
Let $M=\Q(\sqrt{-6})$ and let $E_b/M$ be either cyclic quartic
extension selected in Lemma~\ref{lem:cm-label}.  There is a global unit
$\varepsilon_b\in E_b^\times$ such that, for $g\in\operatorname{Gal}(E_b/M)$,
\[
 \zeta'_{S_M}(0,g)
 =-\log |g(\varepsilon_b)|_{\rm ord},
\qquad
 |z|_{\rm ord}:=\sqrt{z\overline z}.
\tag{15}
\]
Consequently
\[
 L'_{S_M}(0,\psi_b)
 =-\sum_g\overline{\psi_b(g)}
      \log|g(\varepsilon_b)|_{\rm ord}.
\tag{16}
\]
\end{lemma}

\begin{proof}
Stark's proved rank-one theorem for abelian extensions of an imaginary
quadratic base~\cite{Stark1980} applies to $E_b/M$ with the set
\textup{(14)}. Here $|S_M|=4\ge3$, so the unit clause in Stark's
theorem gives a global unit rather than only an $S$-unit. The script
\path{scripts/dimension_eight_cm_unit_lattice.gp} independently
certifies \texttt{bnfcertify=1} and
$e=|\mu(E_b)|=2$ for both fields. Stark's normalized complex-place absolute value is
$|z|_{\rm ord}^2$; hence its factor $1/e$ becomes coefficient one in
the ordinary modulus convention used above.
The finite set contains the distinct conductor primes above $2,3,5$.
Fourier transforming the displayed partial-zeta formula gives
\textup{(16)}.
\end{proof}

Exact unit-action matrices show that the quartic anti-unit component
in Lemma~\ref{lem:stark-normalization} is a rank-two integral lattice.
Rigorous Arb inversion isolates the four coordinate pairs
\[
 (0,2),\quad(-2,0),\quad(0,-2),\quad(2,0)
\tag{17}
\]
with radii below $5\cdot10^{-9}$.  Exact identities in the common
normal closure then match their absolute norms to the original
real-quadratic units.  Thus the two oriented identities required by
the AFK table hold exactly, rather than only in absolute value.

The signed packet and $\Q(\zeta_{16})$ lie in a
convention-selected degree-$128$ compositum.
\path{scripts/dimension_eight_exact_tcc.gp} verifies trace one,
idempotency, and all $784$ minors for both shifts.

\subsection{Conductor one: quadratic units and exact phases}

For $Q=\langle1,-3,1\rangle$, put
\[
 C=\begin{pmatrix}3&-1\\1&0\end{pmatrix},\qquad
 A_8=C^6=\begin{pmatrix}377&-144\\144&-55\end{pmatrix}.
\tag{18}
\]
The one-place and both-place ray groups are $C_2^2$ and $C_2^3$.
Kopp's exponent is one, and the difference is supported on two
quadratic characters.  Writing $\phi=(1+\sqrt5)/2$, their fields and
oriented units are obtained from the positive roots
\[
 h=\sqrt\phi>0,\qquad r=\sqrt{2\phi}>\phi:
\]
\[
\begin{array}{c|c|c}
 &\text{absolute equation}&\text{unit}\\ \hline
L_0&h^4-h^2-1=0&u_0=\phi+h\\
L_1&r^4-2r^2-4=0&u_1=(r+\phi)/(r-\phi).
\end{array}
\tag{19}
\]
Both fields have class number one, with successful unconditional PARI
certificates.

After dropping torsion, their fundamental-unit coordinates are
\[
 \phi=(-2,0),\quad u_0=(0,-1)
 \quad\hbox{in }L_0,\qquad
 \phi=(1,0),\quad u_1=(-1,-2)
 \quad\hbox{in }L_1.
\]
Both determinant indices have absolute value two. This is the exact
relative regulator-index calculation; the analytic class-number
formula therefore gives
\[
 L'(0,\chi_0)=\log u_0,\qquad
 L'(0,\chi_1)=\log u_1.
\tag{20}
\]

\begin{lemma}[Exact magnitude transform]\label{lem:d8-magnitudes}
Let $c=(c_1,c_2)$ be the ray-$8$ class log of a primitive
characteristic.  Its differenced partial-zeta derivative is
\[
 Z'_c(0)=\frac12\left(
 (-1)^{c_2}\log u_0+
 (-1)^{c_1+c_2}\log u_1\right).
\tag{21}
\]
Hence, for
\[
 x^2=u_0,\qquad d^2=\sqrt{u_1/u_0},\qquad a=xd,
\tag{22}
\]
the four ray-$8$ magnitudes are
\[
 [0,0]\mapsto a,\quad [0,1]\mapsto a^{-1},\quad
 [1,0]\mapsto d^{-1},\quad [1,1]\mapsto d.
\tag{23}
\]
At modulus four the stabilizer power is two and
\[
 Z'_c(0)=2(-1)^c\log u_0,
\tag{24}
\]
giving magnitudes $x^2,x^{-2}$; at modulus two the magnitude is one.
These cases exhaust all $63$ nonzero characteristics.
\end{lemma}

\begin{proof}
For the ray group $G=C_2^2$, Fourier inversion gives
\[
 Z'_c(0)=\frac1{|G|}
 \sum_{\chi\in\widehat G}
 \overline{\chi(c)}(1-\overline{\chi(R)})L'(0,\chi).
\]
The sign class is $R=[0,1]$, so only
$\chi_0=[0,1]$ and $\chi_1=[1,1]$ survive, each with coefficient
$1/2$.  Substitution of \textup{(20)} proves \textup{(21)} and taking
half the exponential gives \textup{(23)}.  The identical calculation
for the ray-$4$ group $C_2$, followed by the exactly audited stabilizer
power two, proves \textup{(24)}.  The characteristic audit prints
$12$ occurrences of each ray-$8$ label, six of each ray-$4$ label,
and three modulus-two labels.
\end{proof}

The continued-fraction word of $A_8$ is
$[3,3,3,3,3,3,0]$.  Put $\beta=(3+\sqrt5)/2$.
The exact Euclidean recursion gives the rows
\[
\begin{pmatrix}
377&-144\\144&-55\\55&-21\\21&-8\\
8&-3\\3&-1\\1&0\\0&1
\end{pmatrix},
\tag{25}
\]
whose values on $(\beta,1)^T$ are
\[
\beta^7,\beta^6,\ldots,\beta,1.
\tag{26}
\]
Thus all six adjacent period ratios in AFK's continued-fraction
cocycle are $\beta$.  For the characteristic $(p,q)/8$, division of
the first period expression by the six successive denominators gives
\[
 z_i=\frac{q\beta^{i+2}-p\beta^{i+1}}8,\qquad 0\leq i<6,
\tag{27}
\]
and the outer finite-product index is
\[
 \frac{-144p+(377-1)q}{8}=-18p+47q.
\tag{28}
\]
These identities are checked in exact arithmetic by the archived
maximal-tuple audit.

Put $e(t)=e^{2\pi i t}$ and
\[
 [z;\omega]_n=
 \begin{cases}
 \prod_{k=0}^{n-1}(1-e(z+k\omega)),&n\geq0,\\
 \prod_{k=1}^{-n}(1-e(z-k\omega))^{-1},&n<0.
 \end{cases}
\tag{29}
\]
For $s(z)=\lfloor-z\rfloor+1$, define
\[
 \Sigma(z)=[z/\beta;-1/\beta]_{-s(z)}
 e\!\left(
 \frac{6(z+s(z))^2+6(1-\beta)(z+s(z))}
 {24\beta}\right)
 S_2(z+s(z)+1\mid\beta,1).
\tag{30}
\]
Iterating AFK Definition~1.32 along the rows \textup{(25)}, inserting
its outer finite product \textup{(28)}, and replacing each of Kopp's
double sines by the reciprocal convention \textup{(1)} gives
\[
 \nu_{p,q}=
 \frac{\tau^{-3(p^2-3pq+q^2)}(-1)^{(1+p)(1+q)}}
 {[(q\beta-p)/8;\beta]_{-18p+47q}}
 \prod_{i=0}^{5}
 \Sigma\!\left(\frac{q\beta^{i+2}-p\beta^{i+1}}8\right),
 \quad \tau=-e^{\pi i/8}.
\tag{31}
\]

For real $t\notin\Z$,
\[
 1-e^{2\pi i t}=-2i e^{\pi i t}\sin(\pi t).
\tag{32}
\]
The fundamental reciprocal double sine is positive on
$0<z<\beta+1$, since it is a ratio of positive Barnes double-gamma
values.  Equations
\textup{(2)}, \textup{(31)}, and \textup{(32)} therefore reduce every
phase to an exact element of $\Q(\beta)/2\Z$.
\path{scripts/dimension_eight_maximal_sign_audit.py} proves for all
$63$ nonzero characteristics that the coefficient of $\beta$ cancels,
the remaining coefficient of $\pi$ is integral, and its parity agrees
with the signed radical table.  Decimal arithmetic is used only to
guess a floor; both defining inequalities are then verified exactly in
$\Q(\sqrt5)$.  The audit also proves that no finite-product factor or
recurrence sine vanishes.

With $x,d,a$ as in Lemma~\ref{lem:d8-magnitudes},
the table uses only
\[
 \pm1,\quad\pm x^2,\quad\pm x^{-2},\quad
 \pm d,\quad\pm d^{-1},\quad\pm a,\quad\pm a^{-1}.
\tag{33}
\]
Exact shared-subfield identities glue it to $\Q(\zeta_{16})$.
\path{scripts/dimension_eight_maximal_exact_tcc.py} verifies, for
each shift,
\[
 \Tr\Pi=1,\qquad \Pi^2=\Pi,\qquad
 784\text{ rank-two minors vanish}.
\tag{34}
\]

\subsection{Form-class completion}

Solving the admissibility recurrences gives
\[
 d=8\Longrightarrow
 r=1,\quad n=9,\quad K=\Q(\sqrt5),\quad
 j=2,\quad m=1,\quad f_2=3.
\tag{35}
\]
Hence the form conductor is exactly $f=1$ or $f=3$, and the only
discriminants are $5$ and $45$.  The exact wide class calculations are
\[
 h_{\mathrm{GL}_2}(5)=h_{\mathrm{GL}_2}(45)=1.
\tag{36}
\]
For discriminant five, the maximal order also has narrow class number
one because $\phi$ has norm $-1$.  For discriminant $45$ only the wide
ring class number is used: determinant-$-1$ transformations are
allowed, and no narrow-class-number assertion is needed.

If $t_M$ is obtained from $t$ by $M\in\mathrm{GL}_2(\Z)$, the
unconditional covariance calculation of the companion
Paper I~\cite{ZhaoI} sends
the complete TCC sum at $(t_M,\boldsymbol p)$ to that at
$(t,\ell M\boldsymbol p)$, after reindexing its representatives; for
$\det M=-1$ both cocycle factors are conjugated.  This proves
$\mathcal Z_t\subseteq\mathcal Z_{t_M}$.  Applying the same calculation
to $M^{-1}$ proves the reverse inclusion.  Equations
\textup{(35)}--\textup{(36)} therefore transport each of the two
certified representatives across its own discriminant and complete
dimension eight.

\begin{proof}[Proof of Theorem~\ref{thm:main}]
Section~\ref{sec:d7} proves both shifts for the discriminant-$32$
dimension-seven representative and transports them within that
discriminant. It explicitly excludes the independent discriminant-$8$
stratum. Section~\ref{sec:d8} proves both dimension-eight conductors
and transports within each of the only two possible discriminants.
\end{proof}

\section{Why dimension six remains different}

The dimension-six primitive character has order six.  Its projective
quotient has order three, rather than the order-two quotient that
produces the $V_4$ CM descent in dimension eight.  The alternative
index-two induction groups are nonabelian, so no second abelian
induction base supplies an imaginary-quadratic elliptic-unit theorem.
The finite dimension-six candidate already satisfies all $225$ minors;
the remaining input is one oriented primitive order-six special value.
Thus the gap is analytic, not a finite reconstruction problem.

\section{Reproducibility}

The final audits were run on Linux 6.8.0-124-generic (x86-64, four
virtual AMD EPYC 9354P cores) with PARI/GP 2.15.4,
Python 3.12.3, NumPy 1.26.4, python-flint 0.9.0, and FLINT 3.6.0.
Representative measurements on that virtual machine were
\[
\begin{array}{lrr}
\text{audit}&\text{wall time}&\text{peak RSS}\\ \hline
\text{dimension-seven closure suite}&27.5\ {\rm s}&175\ {\rm MB}\\
\text{Paper-II archive smoke suite}&43.0\ {\rm s}&175\ {\rm MB}.
\end{array}
\]
The submission-specific archive is built deterministically, carries a
self-verifying root manifest named
\texttt{ARCHIVE\_CONTENTS.sha256}, and passes its tests after clean
extraction. Its Zenodo DOI will be inserted after the immutable
deposit; no DOI is claimed in this version.

The supplementary archive contains:
\begin{itemize}
\item all dimension-seven conductor-lowering, Shintani-divisor,
  Arb-height, Artin-label, field-gluing, and exact-minor scripts,
  including the complete divisor and enclosure transcripts;
\item all conductor-three dimension-eight order/maximal-order,
  linear-reinduction, CM-unit, Arb-orientation, normal-closure, and
  exact-minor scripts;
\item all maximal-order dimension-eight ray, quadratic unit, exact
  phase, field gluing, and exact minor scripts;
\item generated CM-orientation and $63$-sign transcripts, the global
  regression transcript, and SHA-256 checksums; and
\item end-to-end regression tests for both dimension-eight strata.
\end{itemize}

The principal reproduction commands are
\begin{verbatim}
python3 -m unittest tests.test_dimension_seven_closure
gp -q scripts/dimension_seven_admissible_strata.gp
gp -q scripts/dimension_seven_exact_tcc.gp
python3 scripts/certify_dimension_seven_double_sine.py --tolerance 1e-10
gp -q scripts/dimension_eight_linear_cm_reinduction.gp
gp -q scripts/dimension_eight_cm_unit_lattice.gp
python3 scripts/certify_dimension_eight_cm_orientation.py
gp -q scripts/dimension_eight_cm_real_unit_bridge.gp
gp -q scripts/dimension_eight_exact_tcc.gp
gp -q scripts/dimension_eight_maximal_tuple_audit.gp
gp -q scripts/dimension_eight_maximal_quadratic_units.gp
python3 scripts/dimension_eight_maximal_sign_audit.py
python3 scripts/dimension_eight_maximal_exact_tcc.py
\end{verbatim}

\section*{Acknowledgments}

The author acknowledges the use of OpenAI's GPT-5.6 Sol model for
mathematical brainstorming, symbolic-computation scaffolding, code
review, and manuscript editing.  All claims are the author's
responsibility and are tied to the exact certificates described above.

\begin{thebibliography}{10}
\bibitem{AFK}
D.~M.~Appleby, S.~T.~Flammia, and G.~S.~Kopp,
\emph{A constructive approach to Zauner's conjecture via the Stark
conjectures}, arXiv:2501.03970.
\bibitem{Kopp}
G.~Kopp, \emph{The Shintani--Faddeev modular cocycle: Stark units from
$q$-Pochhammer ratios}, arXiv:2411.06763.
\bibitem{ZhaoI}
H.~Zhao, \emph{Twisted-convolution identities in dimensions four and
five from Shintani ray units}, companion Paper I, 2026.
\bibitem{Shintani1978}
T.~Shintani, \emph{On certain ray class invariants of real quadratic
fields}, J.\ Math.\ Soc.\ Japan \textbf{30} (1978), 139--167,
\url{https://doi.org/10.2969/jmsj/03010139}.
\bibitem{Voutier1996}
P.~Voutier, \emph{An effective lower bound for the height of algebraic
numbers}, Acta Arith.\ \textbf{74} (1996), 81--95,
\url{https://doi.org/10.4064/aa-74-1-81-95}.
\bibitem{Stark1980}
H.~M.~Stark, \emph{$L$-functions at $s=1$. IV. First derivatives at
$s=0$}, Adv.\ Math.\ \textbf{35} (1980), 197--235,
\url{https://doi.org/10.1016/0001-8708(80)90049-3}.
\bibitem{PARI}
The PARI Group, \emph{PARI/GP version 2.15.4},
\url{https://pari.math.u-bordeaux.fr/}.
\bibitem{pythonflint}
F.~Johansson and contributors, \emph{python-flint 0.9.0},
\url{https://github.com/flintlib/python-flint}.
\end{thebibliography}

\end{document}
